\documentclass[11pt]{article}
\usepackage[margin=1.1in]{geometry}
\usepackage{amsmath,amssymb,amsthm}
\usepackage{booktabs}
\usepackage{lmodern}
\usepackage{microtype}
\usepackage[colorlinks=true,linkcolor=blue,citecolor=blue,urlcolor=blue]{hyperref}

\newtheorem{conjecture}{Conjecture}
\newtheorem{theorem}{Theorem}
\newtheorem{lemma}{Lemma}
\newtheorem{remark}{Remark}
\theoremstyle{definition}
\newtheorem{finding}{Finding}
\newtheorem{question}{Question}

\newcommand{\Li}{\operatorname{Li}}
\newcommand{\e}{\mathrm{e}}
\DeclareMathOperator{\ord}{ord}
\newcommand{\SG}{\mathfrak S}
\newcommand{\CT}{C_3}
\newcommand{\M}{\mathcal M}

\title{Twenty-five conjectures from a local--global random model,\\
with machine verification and adversarial audit}
\author{Working draft\thanks{All constants, counts, and audit
results in this paper are reproducible from the accompanying code
repository: \texttt{run\_all.py} regenerates every number.  Primality
of integers beyond $3.3\times10^{24}$ was tested with Baillie--PSW; all
smaller primality claims are deterministic (fixed-base Miller--Rabin).}}
\date{July 27, 2026}

\begin{document}
\maketitle

\begin{abstract}
We state twenty-five conjectures in elementary and analytic number
theory, each derived from the local--global random model of the primes:
Cram\'er's model corrected by Hardy--Littlewood singular series, with
Bateman--Horn as the organizing framework and Borel--Cantelli accounting
for sparse events.  Each statement is calibrated to the strongest form
the heuristic supports and no stronger.  Every computable constant is
computed from its definition, with local admissibility asserted
programmatically; every statement was tested against exact counts, with
success measured by the \emph{shape} of agreement (predicted constant,
square-root residuals, no drift) rather than by the size of the range
searched.  The statements were then subjected to a three-layer
adversarial audit: an independent literature search for prior art,
scripted refutation attempts a decade beyond the original bounds, and
clean-context re-derivations of constants and counts.  The audit refuted
one conjecture as first stated---the exceptional set for $n=p+k^3$
contains the infinite family of cubes $k^3$ with $3k^2-3k+1$ composite,
an algebraic obstruction the probabilistic accounting missed---and it
appears here restated, with the obstruction promoted to a theorem and
the original error documented; a later round of review caught a
priority failure (a statement we carried as new had been published),
which is likewise documented rather than erased.  We are deliberately
conservative about novelty, and separate two kinds of content.  A
benchmark layer---eleven statements previously stated in essence and
attributed, plus a handful of unrecorded but mechanical
specializations of Bateman--Horn---is offered as calibrated test data
rather than as new mathematics.  The research group carries what we
argue is conceptual novelty: a quantified Chebyshev-type race among
twin-prime pairs modulo $5$ and $8$ whose deterministic term is driven
entirely by the pairs $(q^2-2,\,q^2)$, together with a Goldbach lane
race from the same mechanism (with a verified internal null lane) and
a null-control race in which the mechanism is algebraically impossible
and the predicted drift is zero; a finite-modulus deficit law for
least primes in arithmetic progressions, with a reproducible ordering
anomaly that persists on prime moduli, where smooth-modulus
singular-series effects are absent; derived covariance kernels for the
residual fields of the de Polignac and Hardy--Littlewood-F families
(for the latter the computed cross-member kernel is null, localizing
the observed sub-Poisson profile in the diagonal); a family-uniform
quadratic de Polignac law; an alternating cyclotomic chain and the
complete twin-base analysis of the quadratic cyclotomic family; the
finiteness of factorial twins (the convergent-side Borel--Cantelli
exemplar); a finite-$x$ variance law for prime counts in microscopic
windows with second-order constant $\gamma+\log2\pi-1$ (a microscopic
extrapolation and test of Montgomery--Soundararajan at Gallagher's
boundary); an entanglement-aware local-global law for the mixed
polynomial--exponential sequence $n^2+2^n$; a liminf first-occurrence
law for prime gaps with the Granville-corrected constant
$\sqrt{\e^{\gamma}/2}=0.943682\ldots$, registered against Shanks's
classical constant $1$; and the power-obstruction ladder for
$n=p+j^k$ (total algebraic obstruction for composite $k$, a
Bateman--Horn lane for prime $k$---the dichotomy itself corrected by a
third external review, which refuted our first even/odd version).
\end{abstract}

\tableofcontents

\section{Introduction}

A good conjecture in the theory of numbers is not merely a statement
that has resisted counterexample.  It is the \emph{prediction of a
model}: the assertion that the primes, beyond their local structure,
behave like a random set of density $1/\log n$.  The discipline this
imposes is exacting.  A conjecture worth stating should
(i) survive every congruence and size obstruction, which is a
computation, not a hope;
(ii) come with a quantitative asymptotic whose constant is derived, not
fitted;
(iii) demand exactly the fluctuations the model earns---square-root
noise with logarithmic corrections---and no fewer (the Mertens
conjecture died of demanding fewer);
(iv) sit inside the standard hierarchy of conjectures
(Hardy--Littlewood $k$-tuples $\subset$ Schinzel's Hypothesis~H
$\subset$ Bateman--Horn), so that its failure would propagate; and
(v) be falsifiable instance-by-instance by computation.
Numerical range is the least important column in the ledger: most deep
phenomena drift at the rate $\log\log x$, and $\log\log 10^{18}\approx
3.7$, so verification ``to $10^{18}$'' is weak evidence by itself.  What
persuades is the \emph{shape} of the agreement---a count that tracks a
derived constant through several decades with residuals that look like
noise.

The twenty-five statements of this paper were produced systematically
to those specifications, together with the audit that the
specifications demand.  They are organized in four blocks:
Bateman--Horn systems (\S\ref{sec:bh}), barely-divergent sparse
sequences (\S\ref{sec:sparse}), representation problems with
Borel--Cantelli accounting (\S\ref{sec:rep}), and statistical laws of
the prime sequence (\S\ref{sec:stat}).  Section~\ref{sec:audit}
describes the adversarial audit: what was attacked, what survived, and
in one case what did not.  We consider the audit trail part of the
mathematical content.  Its headline is instructive: the one outright
refutation (Conjecture~\ref{c13}, first version) was produced not by
bad luck but by a density-zero algebraic family---the cubes---invisible
to a probability model that integrates over density.  This is precisely
the classical lesson of admissibility, relearned here in public.

Two further audit outcomes deserve emphasis.  First, an independent
literature layer found that eleven of the twenty-five statements
already exist in essence---some verbatim (Dubner's twin-sum conjecture,
the Stern list for $p+2k^2$, Caldwell--Gallot's $\e^{\gamma}\log N$
laws), one with an exact constant catalogued as an OEIS entry
(A188596, the sub-chain benchmark retained inside
Conjecture~\ref{c6})---and, in one instance caught only at the fourth
round of review, a statement we had carried as new was found already
published (Finding~\ref{f:c10}).  We keep the attributed statements
in the paper, restated in our uniform framework and re-verified at
our bounds, with explicit attribution in \S\ref{sec:attrib}; the
residual claim to novelty of this paper rests on the research group
below, each member of which has now been novelty-searched at the
depth that the Finding~\ref{f:c10} failure taught us to demand.  Second, several of our
\emph{first drafts} were corrected by data before any external
audit: the naive Poisson claim behind Conjecture~\ref{c20} fails at
every accessible scale (the corrected second-order law is the
conjecture); the first-guess bias direction in the twin race of
Conjecture~\ref{c21} was reversed by a mod-$3$ computation; the naive
$\mathrm{Exp}(1)$ mean in Conjecture~\ref{c22} is off by a factor that
took three controls to localize.  We report these corrections because
the difference between a fitted claim and a derived claim is exactly
what an audit is for.

\subsection*{Taxonomy: benchmarks versus research conjectures}

An external review of an earlier draft made a criticism we accept in
full: an unrecorded specialization of Bateman--Horn is analogous to
evaluating a classical special function at a new argument---worthwhile
as data, but not a new conjecture in the conceptual sense.  We
therefore classify the twenty-five statements explicitly.
\emph{Benchmarks} (Conjectures \ref{c2}, \ref{c3}, \ref{c5},
\ref{c7}, \ref{c12}--\ref{c15}, \ref{c17}, \ref{c18}, \ref{c23},
together with the attributed cores of \ref{c1} ($d=2$), \ref{c13},
\ref{c16}(i) and \ref{c24}(i)): statements previously known in
essence, or that follow mechanically from
Hardy--Littlewood/Bateman--Horn once admissibility is checked; their
role here is calibration---each carries a derived constant that the
data must and do hit.  \emph{Research conjectures} (Conjectures
\ref{c1} (the family law), \ref{c4}, \ref{c6}, \ref{c8},
\ref{c9}--\ref{c11}, \ref{c16}(ii), \ref{c19}--\ref{c22},
\ref{c24}(ii), \ref{c25}): statements involving a mechanism, law,
family uniformity, or dichotomy we could not find in the literature,
stated with their dependence assumptions explicit.  This roster is
the product of four successive external reviews, each of which raised
the floor: slots whose occupants were judged bibliographically novel
at most---or, in one case, found to be previously published
(Finding~\ref{f:c10})---were replaced or lifted to family level (the
power-obstruction ladder \ref{c4}, the complete twin-base family
\ref{c5}, the alternating cyclotomic chain \ref{c6}, the
null-mechanism race \ref{c8}, the factorial twins \ref{c10}, the
Goldbach lane race \ref{c25}, and the family lift of \ref{c1});
every replacement was novelty-searched at abstract depth and
verified before admission.  Within the research group we regard the
square-contamination mechanism family---\ref{c21} and \ref{c25} with
their negative control \ref{c8}---together with
Conjecture~\ref{c22}(ii) and the derived covariance kernels of
\ref{c16}(ii) and \ref{c24}(ii) as the paper's core, joined by
\ref{c20} and \ref{c4}; Conjectures \ref{c9}, \ref{c11} and
\ref{c19} carry open modeling questions flagged in place.

\subsection*{Verification standards}

All counts below are exact (sieves; deterministic Miller--Rabin for
$n<3.3\times10^{24}$; Baillie--PSW beyond, flagged as such).  Every
comparison reports the Poisson-normalized residual
$z=(\mathrm{obs}-\mathrm{pred})/\sqrt{\mathrm{pred}}$.  For sparse
counts the cumulative comparison is contaminated at small $x$ by the
phantom mass of the main-term integral near its lower limit; in those
cases we also report per-decade increments, which are clean.  Constants
are computed by truncated Euler products with the truncation wobble
(difference between cutoffs) reported as an error bar; for conditionally
convergent quadratic and cubic products this wobble, not the last digit,
is the honest precision.

\section{Notation}

$p,q$ denote primes; $\gamma$ is Euler's constant,
$\e^{\gamma}=1.781072\ldots$; $\varphi$ is Euler's totient;
$\phi=(1+\sqrt5)/2$; $F_n$ the Fibonacci numbers; $p\#$ the primorial.
For irreducible $f_1,\dots,f_k\in\mathbb Z[x]$ with positive leading
coefficients, the Bateman--Horn singular series is
\[
  C(f_1,\dots,f_k)\;=\;\prod_{p}\frac{1-\omega(p)/p}{(1-1/p)^{k}},
  \qquad
  \omega(p)=\#\{n \bmod p:\ p\mid f_1(n)\cdots f_k(n)\},
\]
and the corresponding main term is
\[
  I(N)\;=\;\int_{2}^{N}\frac{dt}{\prod_{i}\log f_i(t)},
\]
which absorbs the usual $1/\prod_i \deg f_i$.  The Bateman--Horn
conjecture \cite{BH} asserts
$\#\{n\le N: \text{all } f_i(n) \text{ prime}\}\sim C\cdot I(N)$
whenever $C\neq0$; every conjecture in \S\ref{sec:bh} is an instance
with its constant evaluated.  For nonlinear systems the product
converges only conditionally; throughout, it is taken in the canonical
ordering of increasing primes, in which convergence follows from the
equidistribution of $\omega(p)$ about its Chebotarev mean.  Our
numerical truncation wobble is an error bar for this canonically
ordered product, not a substitute for that convergence statement.  We write
$\Li_k(x)=\int_2^x(\log t)^{-k}\,dt$.  The twin singular series is
$\SG(d)=2C_2\prod_{p\mid d,\,p>2}\frac{p-1}{p-2}$ for even $d$, with
$2C_2=1.3203236\ldots$; the Goldbach series $\SG(n)$ is the same product
over odd $p\mid n$.

Admissibility ($\omega(p)<p$ for all $p$) is asserted programmatically
during every constant computation.  The check does real work: the
natural-looking pair $\{n^2+n+1,\ n^2+n+3\}$ has $\omega(3)=3$ and is
rejected; the $k=4$ branch of Conjecture~\ref{c5} and the naive
iterated chain behind Conjecture~\ref{c6} were both rejected by the
same machinery, and both rejections became part of the statements.

Novelty labels, from the independent literature audit
(\S\ref{sec:attrib}): \textbf{(a)} previously stated in essence;
\textbf{(b)} classical family, quantified instance apparently unstated.

\section{Bateman--Horn instances}\label{sec:bh}

The eight conjectures of this block share one hypothesis-generating
move: choose an admissible system of polynomials, compute its singular
series, and state the count.  (Two members go further: \ref{c1}
states a family-uniform law over all even shifts, and \ref{c8} is the
mechanism family's negative control, a race prediction over the
solutions of \ref{c1}'s system.)  The heuristic is uniform, so we
give it once.  Model the primality of the values $f_i(n)$ as independent events
of probability $1/\log f_i(n)$, corrected by the factor
$(1-\omega(p)/p)/(1-1/p)^k$ at each prime $p$ to account for the joint
local behaviour; multiplying the corrections and integrating gives
$C\cdot I(N)$.

\begin{conjecture}[Uniform quadratic de Polignac; family law apparently
unstated, the $d=2$ instance is (a), cf.\ OEIS A080149,
\cite{Carella}]\label{c1}
For even $d\ge2$ let
$\pi^{\mathrm q}_d(N)=\#\{n\le N:\ n^2+1,\ n^2+1+d \text{ both
prime}\}$ and $C(d)=C(x^2+1,\,x^2+1+d)$.  Then, for every fixed $B>0$,
$\pi^{\mathrm q}_d(N)=C(d)\,I_d(N)\,(1+o(1))$ \emph{uniformly} over
even $d\le(\log N)^{B}$.  In particular ($d=2$, the classical instance
that formerly stood alone in this slot):
$\#\{n\le N:\ n^2+1,\ n^2+3 \text{ both prime}\}\sim C(2)\,I(N)$,
$C(2)=2.954014\ldots$
\end{conjecture}

This is the quadratic analogue of the uniform de Polignac law of
Conjecture~\ref{c16}: not one count but the whole profile of constants
must be matched at once, and the family, unlike its linear parent, has
no known prior statement.  Every even shift is admissible (for
$p\ge5$, $\omega(p)\le4<p$; the classes mod $3$ and $5$ vary with $d$,
which is what makes the $C(d)$-profile nonconstant, spanning a factor
${\approx}3$ over $d\le300$).  At $d=2$ the local analysis reads:
$n$ even, $n\not\equiv0\pmod3$,
$\omega(p)=2+\bigl(\frac{-1}{p}\bigr)+\bigl(\frac{-3}{p}\bigr)$ for
$p\ge5$, and the conjecture implies infinitely many twin pairs
$(m+1,m+3)$ with $m$ a perfect square.  Profile verification at
$N=10^6$ over all $150$ even shifts $d\le300$: correlation
$0.99976$ between observed and predicted counts, regression slope
$1.0022$, residual mean $z=+0.24$ with spread $0.84$ and
$\max_d|z|=2.28$---the shape-of-agreement standard applied across the
family, not cherry-picked at one shift.

\begin{conjecture}[Cubic shift; (b)]\label{c2}
$\#\{n\le N:\ n^3+2 \text{ prime}\}\sim C_2^{\ast}\, I(N)$,
$C_2^{\ast}=C(x^3+2)=1.298435\pm0.0003$.
\end{conjecture}

For $p\equiv2\pmod3$ the cube map is a bijection mod $p$, so
$\omega(p)=1$ \emph{exactly} and the local factor is
$(1-1/p)^{-1}(1-1/p)\allowbreak=1$: only $p\equiv1\pmod 3$ (where
$\omega\in\{0,3\}$ according as $-2$ is a cubic residue) moves the
product, which therefore converges only conditionally---hence the quoted
wobble.  No statement of Bunyakovsky type is proven for any cubic; the
nearest theorem is Heath-Brown's $x^3+2y^3$ \cite{HB}, which is the
partial-result neighbour that makes the single-variable statement a
research programme rather than a monolith.

\begin{conjecture}[Chernick chain; (a), Dubner \cite{Dubner2002}]\label{c3}
$\#\{p\le x:\ p,\ 2p-1,\ 3p-2 \text{ all prime}\}\sim C_3\, I(x)$,
$C_3=C(x,2x-1,3x-2)=2.858249\ldots$
\end{conjecture}

The three primes form an arithmetic progression with common difference
$p-1$, first instance $(3,5,7)$; for $p=6m+1$ the product
$p(2p-1)(3p-2)=(6m+1)(12m+1)(18m+1)$ is Chernick's universal-form
Carmichael number, so this conjecture quantifies the population from
which the classical Carmichael constructions draw.

\begin{conjecture}[The power-obstruction ladder: an elementary
structural proposition with a Bateman--Horn corollary; apparently new
as a family]\label{c4}
For $k\ge2$ and $m\ge2$, write $D_k(m)=m^k-(m-1)^k$, and call $m^k$
representable if $m^k=p+j^k$ for some prime $p$ and $j\ge1$.  Then:
\emph{(i)} [theorem] for \emph{composite} $k$, no $k$-th power is
representable: if $r$ is a proper divisor of $k$, then for every
$1\le j<m$ the difference $m^k-j^k$ has the proper factor $m^r-j^r>1$
with cofactor $>1$;
\emph{(ii)} [elementary] for \emph{prime} $k$ (including $k=2$),
$m^k$ is representable if and only if $D_k(m)$ is prime; and $D_k$ is
irreducible for prime $k$ (its roots $\zeta^j/(\zeta^j-1)$, $\zeta$ a
primitive $k$-th root of unity, form a single Galois orbit---$D_k$ is a
linear-fractional transform of the cyclotomic polynomial $\Phi_k$);
\emph{(iii)} [conjecture] for each prime $k$ the representable lane
follows Bateman--Horn:
$\#\{m\le M:\ D_k(m) \text{ prime}\}\sim C_k\int_2^M dt/\log D_k(t)$.
\end{conjecture}

\begin{finding}\label{f:c4}
Our first version of this ladder claimed an even/odd dichotomy, with a
Bateman--Horn lane for every odd $k$.  The third external review
refuted it: for odd composite $k=rs$ the factorization in (i) applies
(e.g.\ $D_9(m)$ is \emph{never} prime), so the lane exists only for
prime $k$.  We verified the refutation directly ($k=9,15,21,25,27,33$:
zero prime values of $D_k$ to $m=2000$, with the predicted divisor
$D_r(m)\mid D_k(m)$ checked identically) and record that this is the
second algebraic-factorization failure the audit programme has caught
in our own output---the exact failure mode our procedure was built to
detect, twice missed at generation time and twice caught downstream.
\end{finding}

The family grew out of the audit refutation behind
Conjecture~\ref{c13} (the case $k=3$).  It replaces, at the suggestion
of an external review, a bibliographically-novel-at-most quintuplet
constellation that formerly occupied this slot (its data survive in the
repository; its $4\cdot10^9$ verification, ratio $1.0004$, remains part
of the audit record as Hardy--Littlewood calibration).

\begin{conjecture}[Twin cyclotomic bases: the complete quadratic
family; no prior trace found]\label{c5}
For the three cyclotomic polynomials of degree $2$ (indices $k$ with
$\varphi(k)=2$, i.e.\ $k\in\{3,4,6\}$), consider
$\Phi_k(n)$ and $\Phi_k(n+1)$ simultaneously prime.  Then:
\emph{(i)} [parity obstruction, elementary] for $k=4$ one member of
the pair $(n^2+1,\,n^2+2n+2)$ is even for \emph{every} $n$, so the
only prime pair is $n=1$, giving $(2,5)$---the same obstruction that
forbids consecutive $n^2+1$ primes;
\emph{(ii)} [collapse, elementary] $\Phi_6(x)=\Phi_3(x-1)$
identically, so the $k=6$ pair at index $n$ \emph{is} the $k=3$ pair
at index $n-1$: the family contains exactly one nontrivial statement;
\emph{(iii)} [conjecture, the live instance] there are infinitely many
$n$ for which the length-3 repunits in bases $n$ and $n+1$ are
simultaneously prime:
$\#\{n\le N:\ \Phi_3(n)=n^2+n+1 \text{ and }
\Phi_3(n+1)=n^2+3n+3 \text{ both prime}\}\sim C_5\,I(N)$,
with $C_5=C(x^2+x+1,\,x^2+3x+3)$.
\end{conjecture}

Both forms in (iii) are odd for every $n$ (each is $\Phi_3$ of an
integer), and mod $3$ each vanishes on exactly one class ($n\equiv1$
and $n\equiv0$ respectively), so the pair is admissible with
$\omega(3)=2$.  Unlike the arbitrary shifted pair that formerly
occupied this slot (replaced at the suggestion of an external review),
the companion here is structurally forced: it is the same cyclotomic
value at the consecutive base, making this the natural ``twin''
statement inside the cyclotomic-chain family of Conjecture~\ref{c6}.
The family analysis (i)--(ii) is what a family \emph{lift} of the
statement uncovers, and both clauses were found by our own
admissibility machinery: the $k=4$ branch was rejected as inadmissible
at $p=2$ at singular-series time (both-odd count to $10^6$: zero,
prime pairs: $n=1$ only), and the $k=6$ run returned counts
\emph{identical} to $k=3$, which is the numerical shadow of the
polynomial identity in (ii).

\begin{conjecture}[The alternating cyclotomic chain; apparently
new]\label{c6}
There are infinitely many primes $p$ such that, writing
$u=\Phi_3(p)=p^2+p+1$, the three integers
\[
p,\qquad u=p^2+p+1,\qquad
\Phi_6(u)=u^2-u+1=p^4+2p^3+2p^2+p+1
\]
are all prime, with
$\#\{p\le x\}\sim C_6^{\mathrm{ch}}\,I(x)$ for the Bateman--Horn
system $\{x,\ x^2+x+1,\ x^4+2x^3+2x^2+x+1\}$,
$C_6^{\mathrm{ch}}=3.6143\pm0.011$.  The first chains are
$(2,7,43)$ and $(3,13,157)$.
\end{conjecture}

This is the repunit analogue of a Cunningham chain: a prime base $p$,
its length-3 repunit $u=(p^3-1)/(p-1)$, then $(u^3+1)/(u+1)$.  The
chain must \emph{alternate} between $\Phi_3$ and $\Phi_6$: the naive
iterate $\{p,\ \Phi_3(p),\ \Phi_3(\Phi_3(p))\}$ is
\emph{inadmissible}---if $p\equiv2\pmod3$ then $u\equiv1\pmod3$, so
$3\mid\Phi_3(u)$ always, while $p\equiv1\pmod 3$ forces $3\mid u$; our admissibility check rejected the naive chain at design time,
and the alternating form is the unique degree-$\{1,2,4\}$ continuation
that survives.  The length-2 sub-chain $\{p,\ p^2+p+1\}$ is the
classical cyclotomic Sophie Germain statement that formerly occupied
this slot ($C=1.521661\pm0.0005$, catalogue value $1.5217315\ldots$,
OEIS A188596; it remains verified at $z=-1.06$ at $10^7$ and is
retained as calibration): infinitely many prime bases have prime
repunits of length $3$, the base-independent home of the phenomenon
whose base-2 shadow is the Mersenne sequence.  The length-3 chain
stated here is, as far as we can find, unstated, and its quartic layer
makes it the highest-degree system in the suite.

\begin{conjecture}[Downward companion; (b); sequence is OEIS
A062326]\label{c7}
$\#\{p\le x:\ p^2-2 \text{ prime}\}\sim C_7\,I(x)$,
$C_7=C(x,x^2-2)=3.383216\ldots$
\end{conjecture}

Squares are $\equiv0,1\pmod3$, so $3\nmid p^2-2$ always---no loss at
$3$---and $\omega(q)=1+\bigl(\frac2q\bigr)$ for the quadratic factor.
The same pairs $(p,\,p^2-2)$ reappear as the deterministic term of the
twin race in Conjecture~\ref{c21}: this conjecture's counting function
sets that term's magnitude.  (The dependence is one-directional and
partial---a race bias could conceivably be established in averaged form
without the full asymptotic here---but a refutation of this conjecture
would remove Conjecture~\ref{c21}'s predicted mechanism.)

\begin{conjecture}[The null-mechanism race; apparently new]\label{c8}
For the quadratic twin pairs of Conjecture~\ref{c1} at $d=2$, square
contamination is \emph{algebraically impossible}: $n^2+1=q^2$ has no
solution with $n\ge1$, and $n^2+3=q^2$ only at $(n,q)=(1,2)$.  The
surviving classes of $n$ mod $5$ are $\{0,1,4\}$, and the classes
$1$ and $4$ have identical singular series (elementary, via CRT).
Let $\pi^{\mathrm q}_a(x)$ count solutions with $n\equiv a\pmod5$ and
$D(x)=\pi^{\mathrm q}_1(x)-\pi^{\mathrm q}_4(x)$.  Then, in contrast
to the twin race of Conjecture~\ref{c21}, this race is
\emph{driftless}: the normalized difference
$D(x)/\sqrt{\pi^{\mathrm q}_1(x)+\pi^{\mathrm q}_4(x)}$ has
logarithmic mean tending to $0$, and no clause analogous to
Conjecture~\ref{c21}(i) holds---there is no deterministic term at the
contamination scale, and leadership wanders in arcsine fashion with no
persistent leader.
\end{conjecture}

This is the negative control that completes the mechanism family of
Conjectures~\ref{c21} and~\ref{c25}: those races carry a predicted
drift \emph{because} prime squares can infiltrate specific residue
classes of specific patterns; here the quadratic form blocks squares
identically, so the same model predicts \emph{nothing}---a falsifiable
contrast.  If a persistent drift were found in this race, the
mechanism story behind the family would be wrong no matter how well
the positive cases fit.  Verification at $10^7$: $32{,}898$ pairs,
class counts $10{,}917$ against $10{,}981$
($D=-64$, i.e.\ $-0.43$ noise units), log-mean normalized drift
$-0.46$ (consistent with $0$ at one noise unit), weighted lead
fraction $0.21$---deep in arcsine territory, exactly the profile of a
fair race.  (The quadratic triple $n^2+\{1,3,7\}$ that formerly
occupied this slot, $C=10.64599\pm0.006$, is retired to the audit
record with its calibration intact---observed $3{,}963$ against
predicted $3{,}998.6$ at $10^7$, $z=-0.56$, independently recomputed
to $3\times10^7$; the third external review found the longer pattern
$n^2+\{1,3,7,9,13\}$ recorded in the OEIS since 2000, so the tuple
family has prior presence and only its singular-series evaluation was
plausibly new.)

\subsection*{Verification}

\begin{center}\small
\begin{tabular}{lrrrrr}
\toprule
 & bound & observed & predicted & ratio & $z$ \\
\midrule
C\ref{c1}$_{d=2}$  & $10^{7}$ & $32{,}898$ & $32{,}862.7$ & $1.0011$ & $+0.19$ \\
C\ref{c2}  & $10^{7}$ & $287{,}956$ & $287{,}784.8$ & $1.0006$ & $+0.32$ \\
C\ref{c3}  & $3\cdot10^{8}$ & $125{,}379$ & $125{,}429.4$ & $0.9996$ & $-0.14$ \\
C\ref{c4}$_{k=2}$ & $10^{7}$ & $1{,}270{,}606$ & $1{,}270{,}902.8$ & $0.9998$ & $-0.26$ \\
C\ref{c4}$_{k=3}$ & $10^{6}$ & $126{,}826$ & $126{,}641.6$ & $1.0015$ & $+0.52$ \\
C\ref{c4}$_{k=5}$ & $8\cdot10^{5}$ & $56{,}925$ & $57{,}021.3$ & $0.9983$ & $-0.40$ \\
C\ref{c5}  & $10^{7}$ & $33{,}274$ & $32{,}976.7$ & $1.0090$ & $+1.64$ \\
C\ref{c6}  & $10^{7}$ & $1{,}362$ & $1{,}357.2$ & $1.0035$ & $+0.13$ \\
C\ref{c7}  & $10^{7}$ & $74{,}914$ & $75{,}273.8$ & $0.9952$ & $-1.31$ \\
\bottomrule
\end{tabular}
\end{center}

\smallskip
For C\ref{c1}, the family-profile verification ($150$ shifts) is
reported at the statement; the table row is the $d=2$ instance.
For C\ref{c4}: the even-rung theorem was checked numerically to $10^8$
(no fourth power is $p+j^4$); the three verifiable Bateman--Horn lanes
$k=2,3,5$ (constants $2$ exactly, $3.36181$, and $3.67770\pm0.007$, the
quartic product computed by brute-force root counts to $10^5$) all sit
within $|z|<0.6$.  For C\ref{c5}, $C_5=2.964239\pm0.002$.
For C\ref{c6}, the chain census by decade is
$4,\ 11,\ 41,\ 224,\ 1362$ at $N=10^3,\dots,10^7$ ($z$ from
$+0.40$ to $+0.13$, no drift), the quartic layer's singular series
computed by brute-force root counts to $10^5$; the retained
sub-chain benchmark $\{p,\,p^2+p+1\}$ stands at $33{,}661$ against
$33{,}855.5$ at $10^7$ ($z=-1.06$).  For C\ref{c8}, the race data are
reported at the statement.
Independent clean-context recomputation gave
$C_7=3.383227$ and, for the retired occupants of
slots~\ref{c6} and~\ref{c8} (the Sophie Germain sub-chain and the
quadratic triple), $1.521730$ and $10.647706$, with primes to $10^9$
(each inside our stated wobble) and independent counts to $3\times10^7$
with final ratios $0.9962$, $0.9965$, $1.0059$.  (The quintuplet that
formerly occupied slot~\ref{c4} retains its audit-record calibration:
observed $15{,}236$ against predicted $15{,}230.3$ at $4\cdot10^9$,
ratio $1.0004$, $z=+0.05$.)

\section{Sparse sequences: barely-divergent Borel--Cantelli}
\label{sec:sparse}

For sequences whose $n$th term has size $\e^{cn}$ the expected number of
primes is $\sum_n \kappa_n/(c\,n)$: convergence or divergence of this
sum is the entire question, and when it diverges it diverges
logarithmically, so the right conjecture has counting function
$\alpha\log x$ with a derived $\alpha$---never more.

\begin{conjecture}[Fibonacci primes; cf.\
\cite{GranthamGranville}]\label{c9}
$\#\{p\le x:\ F_p \text{ prime}\}\sim c_F\log x$ for a constant
$c_F>0$ determined by the local (rank-of-apparition) structure of the
Fibonacci sequence in the sense of the Grantham--Granville
recurrence-sequence model.  The first-order screening value is
$c_F^{(0)}=\e^{\gamma}/\log\phi=3.70122\ldots$; the data below
indicate $c_F<c_F^{(0)}$, and we deliberately do \emph{not} conjecture
the exact constant.
\end{conjecture}

Every prime factor of $F_p$ ($p\neq5$ prime) has rank of apparition
exactly $p$, hence is $\equiv\pm1\pmod p$: the
Lenstra--Pomerance--Wagstaff screening \cite{Wagstaffheu} transfers with
$\log2\mapsto\log\phi$.  The data sit persistently \emph{below} the
constant: $25$ Fibonacci prime indices $\le10^4$ (independently
reconfirmed by a second implementation) against a predicted $34.1$,
ratio ${\approx}0.73$, $z=-1.6$, stable over three decades.  Grantham
and Granville \cite{GranthamGranville} caution that naive constants for
such sequences require $O(1)$ local corrections; the mod-$5$ divisor
structure of $F_p$ is the natural suspect, and we flag this constant as
the component of the suite most likely to need a Granville-type patch.

\begin{conjecture}[Factorial twins; apparently new]\label{c10}
Only finitely many $n$ have $n!-1$ and $n!+1$ simultaneously prime,
and the complete list is $n=3$: i.e.\ $6=3!=3\#$ is the unique
factorial lying between twin primes.
\end{conjecture}

Each event separately has Caldwell--Gallot probability
$\sim\e^{\gamma}/n$ \cite{CaldwellGallot} and divergent sum (the
$n!+1$ side is exactly Conjecture~\ref{c12}); the \emph{joint} event
has probability $\sim(\e^{\gamma}/n)^2$, whose sum converges.  This is
the suite's exemplar of the convergent side of the Borel--Cantelli
dichotomy---the accounting that predicts finitude rather than
infinitude.  Verified for $n\le700$; the model-expected number of
further examples beyond $700$ is $\e^{2\gamma}/700\approx
4.5\times10^{-3}$ (computed tail).  The defensible core is finiteness;
the exact one-element list is the Goldilocks-maximal form---strictly
stronger than the tail estimate can guarantee, and chosen deliberately
as the most falsifiable statement.

\begin{finding}\label{f:c10}
The statement that previously occupied this slot---primorial twins:
$p\#\pm1$ both prime only for $p\in\{3,5,11\}$---was carried for three
revisions under an ``apparently new'' label, and the fourth external
review found it stated in Lillie \cite{Lillie}, whose abstract gives
both the $O(n^{-2})$ joint probability and the prediction that there
are in total three instances.  The failure is instructive and we
record it against ourselves: our novelty search had \emph{surfaced
that exact preprint}, but the search snippet described only its
single-sided results, and we did not read to the abstract's second
half.  Priority checks on a conjecture must read the abstract and body
of near-neighbour papers, not their search summaries; the audit
protocol now requires it.  The primorial statement remains, attributed
to Lillie, as calibration (our independent verification to $p\le4000$
stands), and the slot now carries the factorial analogue, for which
our deeper search---including the factorial-prime literature and the
Lillie preprint itself---found no prior statement.
\end{finding}

\begin{conjecture}[{$n^2+2^n$}; sequence is OEIS A064539]\label{c11}
Every prime of the form $n^2+2^n$ with $n>1$ has $n\equiv3\pmod6$
(elementary), and:
\emph{(i)} infinitely many such primes exist, with counting function
$\asymp\log N$;
\emph{(ii)} [candidate law, entanglement-aware form] for a finite set
$S$ of primes let $D_S$ be the \emph{exact} joint density, within the
lane $n\equiv3\ (6)$, of $n$ with $p\nmid n^2+2^n$ for all $p\in S$,
computed over the full joint period
$\mathrm{lcm}_{p\in S}\,\mathrm{lcm}(6,\,p\cdot\ord_p2)$, and set
$\kappa_S=D_S/\prod_{p\in S}(1-1/p)$.  The conjecture is that the net
$\{\kappa_S\}$ converges as $S$ increases to all primes, with limit
$\kappa$, and that the count is
$\sim\sum_{n\le N,\,n\equiv3\,(6)}\kappa/\log(n^2+2^n)$.  The working
value, from the CRT-exact core $p\le19$ times per-prime factors to
$p\le300$, is $\kappa=4.2734\ldots$
\end{conjecture}

This is not a Bateman--Horn problem: the local conditions at different
primes are tied through the orders $\ord_p2$, and the per-prime
product $\prod_p(1-\delta_p)/(1-1/p)$ of an earlier draft tacitly
assumed their independence.  Following external review, the law is now
stated through the CRT-exact quantities $\kappa_S$, which make no
independence assumption.  The computation then returned a small
surprise in the other direction: through $p\le19$ (joint period
$116{,}396{,}280$) the exact joint density \emph{factorizes to machine
precision}---$\kappa_S$ agrees with the naive per-prime product to
$10^{-15}$ at every level $S=\{p\le P\}$, $P\in\{5,7,11,13,17,19\}$.
The entanglement the reviewer warned about is thus computably absent
at all accessible levels; whether exact factorization persists for
all $p$ (which would itself be a striking equidistribution statement
about the orbits of $2$) is part of what clause (ii) asks.  We
separate the robust claim (i) from the candidate law (ii), and note
that the eight hits below $6000$ are far too few to validate a
four-digit constant---the verification below checks consistency, not
the constant.

The congruence claim is elementary: $n$ odd is forced by parity, and for
odd $n$ with $3\nmid n$ we have $n^2+2^n\equiv1+2\equiv0\pmod3$.  What
appears to be a harsh obstruction is, inside the surviving lane, a local
bonus: $\delta_2=\delta_3=0$ contribute the factors $2\cdot\tfrac32=3$
to $\kappa$.  Mixed polynomial--exponential sequences are exactly where
hand-waving fails and the computed local product earns its keep.

\begin{conjecture}[Factorial primes; (a),
\cite{CaldwellGallot}]\label{c12}
$\#\{n\le N:\ n!+1 \text{ prime}\}\sim \e^{\gamma}\log N$.
\end{conjecture}

\subsection*{Verification}

C\ref{c9}: hits $\{3,5,7,\dots,9311,9677\}$, $25$ of predicted $34.1$
($z=-1.56$; see the remark above).
C\ref{c10}: exhaustive to $n\le700$: the only factorial twin is
$n=3$; individually, $n!+1$ is prime for
$n\in\{2,3,11,27,37,41,73,77,116,154,320,340,399,427\}$ and $n!-1$
for $n\in\{3,4,6,7,12,14,30,32,33,38,94,166,324,379,469,546\}$---two
divergent single-sided laws straddling their shared constant while
the joint count stops at one, the convergent accounting in action.
(The primorial companion, now attributed to Lillie \cite{Lillie} per
Finding~\ref{f:c10}: exhaustive to $p\le4000$, twins exactly
$\{3,5,11\}$, with $p\#+1$ prime for eleven $p$ against the
Caldwell--Gallot prediction $14.8$, $z=-0.98$.)
C\ref{c11}: hits $\{3,9,15,21,33,2007,2127,3759\}$ up to $6000$:
observed $8$, model $8.55$, $z=-0.19$ (at the earlier bound $4200$:
$8$ vs $8.19$, $z=-0.07$); the lane constraint was verified
exhaustively (no prime off $n\equiv3\ (6)$ to $n=3000$, independently
reconfirmed), and the CRT-exact factorization of clause (ii) was
verified through $p\le19$ as described at the statement.
C\ref{c12}: $15$ primes with $n\le700$ against $11.7$ ($z=+0.98$),
completing the $\e^{\gamma}\log$ family benchmark that
Conjecture~\ref{c10} joins on its convergent side.

\section{Representation problems}\label{sec:rep}

\begin{theorem}[Cube obstruction; found by the adversarial
audit]\label{t:cube}
For $k\ge2$, the cube $k^3$ is representable as $p+j^3$ with $p$ prime,
$j\ge1$, if and only if $3k^2-3k+1$ is prime.  Consequently the set of
integers not representable as $p+k^3$ is infinite: the non-representable
\emph{cubes} alone have counting function $\sim x^{1/3}$.  (That the
full exceptional set is $\sim x^{1/3}$ follows only when combined with
Conjecture~\ref{c13}(i), and is conjectural.)
\end{theorem}

\begin{proof}
$k^3-j^3=(k-j)(k^2+kj+j^2)$; for $1\le j\le k-2$ both factors exceed
$1$, so the difference is composite.  The only candidate is $j=k-1$,
giving $3k^2-3k+1$.  Since $3k^2-3k+1$ is composite for a density-one
set of $k$ (all but $O(K/\log K)$ of $k\le K$, by an upper-bound sieve
of Brun or Selberg type applied to the quadratic's prime values), all
but a vanishing proportion of cubes are unrepresentable.
\end{proof}

\begin{finding}\label{f:c13}
As first generated (``the exception set of $n=p+k^3$, $k\ge1$, is
finite''), the next conjecture was \textbf{false}, and our own
adversarial battery refuted it: all $412$ exceptions found in
$(10^8,10^9]$ are perfect cubes, explained by Theorem~\ref{t:cube}.
The error is a textbook violation of the admissibility clause of the
design criteria---an algebraic obstruction along a density-zero
subsequence, invisible to a Borel--Cantelli sum taken over all $n$.  We
record the failure rather than hide it; the restated version separates
theorem from conjecture.
\end{finding}

\begin{conjecture}[Prime plus a cube, restated; (a)-core:\ the non-cube
finiteness is Hardy--Littlewood's $E_3(X)=O(1)$ \cite{HL1923}]\label{c13}
\emph{(i)} The set of non-cube integers $n\ge2$ not representable as
$p+k^3$ ($k\ge1$) is finite.
\emph{(ii)} The representable cubes obey Bateman--Horn:
$\#\{k\le K:\ 3k^2-3k+1 \text{ prime}\}\sim C(3x^2-3x+1)\,I(K)$.
\end{conjecture}

(Empirically the non-cube exception census by decade is
$4,\,27,\,168,\,763,\,2011,\,2808,\,1181,\,88$ up to $10^8$,
peaking at the sixth decade and collapsing thereafter, with largest
element found $78{,}526{,}384$; these are verification data, not part
of the formal statement.)

\begin{conjecture}[Stern list; (a): OEIS A060003]\label{c14}
Exactly ten odd integers are not of the form $p+2k^2$ with $k\ge1$:
\[
1,\ 3,\ 17,\ 137,\ 227,\ 977,\ 1187,\ 1493,\ 5777,\ 5993.
\]
\end{conjecture}

The expected number of representations of odd $n$ is
$\asymp\sqrt{n}/\log n$, so the expected number of exceptions beyond any
$X$ decays like $\sum_{n>X}\exp(-c\sqrt n/\log n)$: the list should be
complete, and the seven prime members are precisely the Stern primes.
Verified here to $10^9$, one decade past our original bound, with no
eleventh exception.

\begin{conjecture}[Goldbach in the $(3,3)$ lane; (a), K.~Martin
\cite{Martin}]\label{c15}
Every $n\equiv2\pmod4$ with $n\ge6$ is $p+q$ with
$p\equiv q\equiv3\pmod4$, and the ordered representation count obeys
\[
R_3(n)\;\sim\;\tfrac12\,\SG(n)
\int_{3}^{n-3}\frac{dt}{\log t\,\log(n-t)} .
\]
\end{conjecture}

Normalization: $R_3(n)$ counts \emph{ordered} pairs $(p,q)$, and the
full-interval integral is the ordered main term for all Goldbach
representations, so the factor $\tfrac12$ is exactly the
equidistribution split between the $(1,1)$ and $(3,3)$ lanes (parity
forces $p\equiv q\pmod4$ when $n\equiv2\pmod4$).  The $(3,3)$ lane
keeps its smallest representations ($3+q$) and is the lane with
\emph{no} exceptions at all, which is why the uniform threshold
$n\ge6$ can be conjectured.  Verified exhaustively to
$10^9$; the count formula matches to $0.2$--$1.3\%$ at
$n\approx10^6,10^7,10^8$.

\begin{conjecture}[Uniform quantitative de Polignac;
(a)-core]\label{c16}
\emph{(i)} $\pi_d(x)=\#\{p\le x-d:\ p,\ p+d \text{ prime}\}
=\SG(d)\,\Li_2(x)\,(1+o(1))$ uniformly for even $d\le(\log x)^{A}$.
\emph{(ii)} [residual law, with derived covariance kernel]  With
$z_d(x)=(\pi_d(x)-\SG(d)\Li_2(x))/\sqrt{\SG(d)\Li_2(x)}$: the same
primes participate in many $\pi_d$, and the covariance of two pair
counts is carried by triple configurations sharing one prime.  For
even $d\neq d'$ define the overlap constant
\[
K(d,d')=\CT(0,d,d')+\CT(0,d',d+d')+\CT(0,d,d+d')
+\CT\bigl(0,\,|d-d'|,\,\max(d,d')\bigr),
\]
where $\CT$ denotes the Hardy--Littlewood triple constant of the
offset set (degenerate or inadmissible triples contributing $0$).
Then
$\operatorname{Cov}(\pi_d,\pi_{d'})=K(d,d')\,\Li_3(x)\,(1+o(1))$,
so the correlation kernel is
$\rho(d,d')\sim K(d,d')\,\Li_3(x)\big/\bigl(\SG(d)\SG(d')\bigr)^{1/2}
\Li_2(x)$, and the conjecture is that the residual field
$\{z_d(x)\}_{2\mid d\le 2D(x)}$, $D(x)=(\log x)^{A}$, converges to a
centered Gaussian field with \emph{this} covariance, with no trend in
$d$.  The diagonal ($4$-tuple) correction to
$\operatorname{Var}(\pi_d)$ below its Poisson value---the pair
analogue of the Montgomery--Soundararajan deficit of
Conjecture~\ref{c20}---is the remaining underived component, and we
state it as open rather than fold it into a fitted $\sigma$.
\end{conjecture}

Numerical evaluation of the kernel (all even $d,d'\le60$ at
$x=10^8$, triple constants by Euler products to $5\times10^4$): mean
off-diagonal $\rho=0.43$, hence predicted common-mode variance
$\operatorname{Var}(\bar z)\approx0.41$ and predicted profile spread
$\sigma^2\le1-0.41=0.59$ \emph{before} the diagonal correction; the
observed profile variance ${\approx}0.27$ at $x=10^8$ then implies a
diagonal (sub-Poisson) factor ${\approx}0.5$, of the same sign and
order as the Montgomery--Soundararajan mechanism supplies for single
primes.  An earlier draft asserted only ``$\sigma^2\le1$,
determined by the covariance structure''; the external review rightly
called that a task rather than a law, and the kernel above is the
law the task produces.

The force of the statement is its uniformity: the singular series takes
a thousand different values over $d\le2000$, spanning a factor of about
three between $d=2$ and the primorial-rich $d$, and the counts must
reproduce the entire profile simultaneously with square-root errors.  A
conspiracy would need to bend a thousand-dimensional vector, not one
number.  At $x=10^8$: correlation $0.999999$, slope $0.99966$,
$\max_d|z_d|=1.80$; stressed to $d\le6000$, $\max|z_d|=2.01$ over three
thousand values.

\begin{conjecture}[Dubner's conjecture; (a) \cite{Dubner2000}]\label{c17}
Every even $n\ge4210$ is a sum of two members of twin-prime pairs; the
$35$ exceptions (OEIS A007534, largest $4208$) are complete.
\end{conjecture}

Twin members have density $\asymp1/\log^2x$, so expected representations
grow like $n/\log^4 n$ and exceptions must be finite; the observed
exceptions cluster in triples of consecutive even numbers, a structural
fingerprint of the sparse low-lying twin population.  Verified to
$10^8$, adversarially re-swept to $10^9$, and independently re-verified
to $10^9$ by a second implementation with a different algorithm: the
same $35$ exceptions each time.

\section{Statistical laws}\label{sec:stat}

\begin{conjecture}[Twin-gap records; (a), Kourbatov
\cite{Kourbatov}]\label{c18}
\emph{(i)} The maximal gap $G_t(x)$ between consecutive twin-prime
starts up to $x$ satisfies $G_t(x)\asymp\log^3x$.
\emph{(ii)} [candidate constant]
$\limsup_x G_t(x)/\log^3 x = 1/(2C_2)=0.757390\ldots$
Part (i) is the robust claim; the exact constant in (ii) inherits the
known fragility of Cram\'er-style extreme-value constants (Granville's
correction arose precisely because extreme events feel local sieve
scales invisible to an inhomogeneous Poisson model), and is offered as
a working value, not an assertion at the same confidence as (i).
\end{conjecture}

Twin starts thin the integers to density $2C_2/\log^2t$; a Cram\'er
process at that density has record waiting times
$\sim(\text{mean spacing})\times\log(\text{number of events})
=\log^3x/2C_2$.  The Cram\'er--Granville episode teaches that such
constants can be shifted by local corrections (Granville's
$2\e^{-\gamma}$ \cite{Granville}); the claim to take seriously is the
cube of the logarithm, and the constant is flagged as first-order.
Records to $4\cdot10^9$ (largest observed twin-to-twin gap $5292$,
following the twin at $2{,}466{,}641{,}069$) wander in
$[0.41,0.56]\cdot\log^3x$, approaching the asymptote from below on a
$\log\log$ clock, as such statistics must.

\begin{conjecture}[First occurrences of prime gaps, liminf form;
(b)]\label{c19}
Let $\mathcal G$ be the set of even $g$ that occur as gaps between
consecutive primes (all even $g$, under Polignac's conjecture), and for
$g\in\mathcal G$ let $p(g)$ be the prime starting the first such gap.
Then:
\emph{(i)} $\log p(g)/\sqrt g$ is bounded between positive constants
on $\mathcal G$;
\emph{(ii)} [liminf law]
\[
\liminf_{\substack{g\to\infty\\ g\in\mathcal G}}
\ \frac{\log p(g)}{\sqrt g}
\;=\;\sqrt{\e^{\gamma}/2}\;=\;0.943682\ldots,
\]
the liminf being attained along gap sizes realized at the
Cram\'er--Granville maximal-gap envelope.  (The restriction to
$\mathcal G$ repairs a definitional defect caught by external review:
$p(g)$ is not known to be defined for every even $g$, and quantifying
over all $g$ would smuggle Polignac's conjecture into the statement.)
\end{conjecture}

An earlier draft asserted the full limit, inverting Granville's
$\limsup G(x)/\log^2x=2\e^{-\gamma}$; an external review correctly
observed that this inversion is invalid.  A limsup envelope for maximal
gaps controls only where gaps \emph{can first appear at the earliest},
i.e.\ a lower envelope for $p(g)$---hence a liminf claim, not a limit.
Whether $\log p(g)/\sqrt g$ converges at all, and if so whether the
limit is Shanks's classical $1$ \cite{Shanks} (as Wolf's refinement
$p(g)\sim\sqrt g\,\e^{\sqrt g}$ and the Kourbatov--Wolf residue-class
laws \cite{KW2020} also assert) or the envelope value, involves at
least three distinct phenomena---whether a given $g$ occurs near the
extreme scale at all, the distribution of exact gap lengths near
records, and the difference between record gaps and non-record first
occurrences---and is left open here.  The two candidate constants
differ by $6\%$, beneath the resolution of any feasible computation:
our regression slope over $g\ge100$ falls from $1.307$ at $10^9$ to
$1.297$ at $4\cdot10^9$, far above both and decreasing, as it must be
for either.

\begin{conjecture}[Microscopic variance law; (b); the constant is
Montgomery--Soundararajan's \cite{MS}]\label{c20}
Fix $\lambda>0$ and $c\in(0,1]$, and let $X$ count primes in
$[t,\,t+\lambda\log x)$ for $t$ uniform on $[x,(1+c)x]$.  Then, under
the Hardy--Littlewood $k$-tuple conjectures,
$X\Rightarrow\mathrm{Poisson}(\lambda)$ (Gallagher's argument
\cite{Gallagher}, which is conditional on exactly those conjectures),
and at finite $x$
\[
\frac{\operatorname{Var}X}{\mathbb E\,X}
\;=\;1-\frac{\log(\lambda\log x)+\gamma+\log2\pi-1}{\log x}
+o\!\Bigl(\frac1{\log x}\Bigr),
\]
with $\gamma+\log2\pi-1=1.41509\ldots$, the leading correction being
independent of the sampling constant $c$.
\end{conjecture}

The deficit arises by inserting Montgomery's singular-series average
$\sum_{d\le h}\SG(d)(h-d)=\tfrac{h^2}2-\tfrac h2(\log h+\gamma+\log2\pi
-1)+o(h)$ into the second factorial moment at the microscopic window
$h=\lambda\log x$.  We state plainly, following external review, what this is and is not:
the constant $\gamma+\log2\pi-1$ is Montgomery--Soundararajan's, was
identified in their framework long before this paper, and the statement
is a \emph{microscopic extrapolation and numerical test} of their
second-moment formula at Gallagher's boundary---not a newly discovered
second-order law.  The hypothesis must also be quantitative: what is
required is not the qualitative $k$-tuple conjectures but uniform
Hardy--Littlewood estimates for all shifts up to $O(\log x)$, control
of the average singular series at the $o(h)$ level, and bounds for
prime powers, endpoint weights, and the variation of $\log t$ across
the sampling range.  Three further cautions are built into the
statement: Gallagher's Poisson law is conditional, not a theorem about
primes; the sampling measure for $t$ is specified, since different
weightings shift lower-order terms; and the error is claimed only as
$o(1/\log x)$.  The naive claim $\operatorname{Var}/\mathbb E=1$ fails at
every accessible scale by $18$--$28\%$; the corrected law matches with
no fitted parameter, and the observed residual $\approx+0.003$ is
consistent with (but not evidence for) a genuine $1/\log^2x$ term.
(Numerical precedent: Sanchis-Lozano \cite{SL} fitted the
Montgomery--Soundararajan variance shape at $x\sim10^8$ in the
mesoscopic regime $h\gg\log N$; the microscopic statement and test
appear to be new, and we claim no more than that.)

\begin{center}\small
\begin{tabular}{lcccc}
\toprule
 & \multicolumn{2}{c}{$x=10^9$} & \multicolumn{2}{c}{$x=4\cdot10^9$}\\
$\lambda$ & observed & predicted & observed & predicted\\
\midrule
$1/2$ & $0.8241$ & $0.8189$ & --- & ---\\
$1$   & $0.7910$ & $0.7854$ & $0.7998$ & $0.7960$\\
$2$   & $0.7527$ & $0.7520$ & $0.7670$ & $0.7646$\\
$4$   & $0.7201$ & $0.7185$ & --- & ---\\
\bottomrule
\end{tabular}
\end{center}

\noindent The residual $\approx+0.003$ is the anticipated
$O(1/\log^2x)$ term.

Following the fourth external review, the twin-race statement is now
split into its provable algebraic core and its conjectural clauses,
with the fluctuation functional defined rather than gestured at.  For
a locally integrable $F$ write
$\M_x(F)=\frac1{\log x}\int_2^x F(t)\,\frac{dt}t$
for the logarithmic mean.

\begin{lemma}[Orientation elimination and class assignment]\label{l:c21}
In $\psi_2(x)=\sum_{n\le x}\Lambda(n)\Lambda(n+2)$, the total
contribution of terms in which $n$ or $n+2$ is a proper prime power
is, apart from $O(x^{1/3}\log^2 x)$ from cubes and higher powers,
carried entirely by the patterns $(q^2-2,\,q^2)$ with $q$ prime and
$q^2-2$ prime: the mirror pattern $(q^2,\,q^2+2)$ is annihilated by
$3\mid q^2+2$ for every prime $q>3$.  Moreover, for every prime
$q\neq5$, $q^2-2\equiv2$ or $4\pmod 5$ (never $1$), and for every odd
$q$, $q^2-2\equiv7\pmod8$.  (The single exceptional term $q=5$,
$(23,25)$, lands outside the twin-start classes altogether and is
absorbed in the bounded error.)
\end{lemma}

The proof is elementary (squares mod $3$, $5$, $8$; the prime-power
count).  The lemma orients the race: prime-square contamination can
enter only specific residue classes, and the model converts that into
a drift prediction for the \emph{unweighted} twin counts.

\begin{conjecture}[Twin race modulo 5 and 8; (b), apparently
new]\label{c21}
Twin starts $p>5$ lie in classes $p\equiv1,2,4\pmod5$.  Let
$D_1(x)=\pi_t(x;5,1)-\tfrac12\bigl(\pi_t(x;5,2)+\pi_t(x;5,4)\bigr)$
and
\[
T(x)\;=\;\frac1{2\log^2x}
\sum_{\substack{q\le\sqrt x\\ q^2-2\ \mathrm{prime}}}
\log q\,\log(q^2-2),
\]
whose scale is governed by the constant $C_7$ of
Conjecture~\ref{c7}.  Then:
\emph{(i)} [drift law] class $1$ carries the systematic surplus $T$:
\[
\M_x\!\left(\frac{D_1-T}{\sqrt{\pi_t}}\right)\longrightarrow0
\qquad(x\to\infty),
\]
i.e.\ the deterministic part of $D_1$ is exactly $T$, and the
remainder is mean-zero noise at the random-walk scale under
logarithmic averaging;
\emph{(ii)} [null covariance] the symmetric classes carry no
deterministic term:
$\M_x\bigl((\pi_t(\cdot;5,2)-\pi_t(\cdot;5,4))/\sqrt{\pi_t}\bigr)\to0$,
and likewise for each of the three pairwise differences among classes
$\{1,3,5\}$ mod $8$ in clause (iv);
\emph{(iii)} [sign density, conditional] conditional on the Gaussian
fluctuation model for the normalized remainder in (i)--(ii), the
logarithmic density of $\{D_1>0\}$ exists and equals $\tfrac12$,
approached from above with a finite-$x$ excess of order $1/\log x$
(the drift-to-noise ratio $T/\sqrt{\pi_t}\asymp1/\log x$): in
contrast to Chebyshev's races, the twin race has no persistent
leader;
\emph{(iv)} [mod-8 companion] since $q^2-2\equiv7\pmod8$ always
(Lemma~\ref{l:c21}), the \emph{entire} deterministic term falls on
the single class $7\pmod 8$: with
$D_7(x)=\tfrac13\sum_{a\in\{1,3,5\}}\pi_t(x;8,a)-\pi_t(x;8,7)$,
\[
\M_x\!\left(\frac{D_7-2T}{\sqrt{\pi_t}}\right)\longrightarrow0
\]
(twice the mod-5 term, undiluted), while classes $1,3,5$ are mutually
symmetric.  The mod-5 and mod-8 races are two projections of one
mechanism, so their joint behaviour is a family test in the sense of
the shape-of-agreement criterion.
\end{conjecture}

An earlier draft of clause (iii) asserted a limiting sign density
strictly between $\tfrac12$ and $1$; an external review correctly
observed that this is inconsistent with the scale analysis---when
drift/noise $\to0$, a Gaussian model forces the sign density to
$\tfrac12$, and a nontrivial limiting density would require a
persistent normalized bias that no mechanism here produces.  A yet
earlier draft wrote the error in (iv) as
$O^{*}(\sqrt{\pi_t})$; the fourth review correctly objected that
$O$-notation asserts a provable bound where only a conjectural
fluctuation statement is meant, and the $\M_x$-clauses above say
exactly what is conjectured and no more.  Two model assumptions
remain open and are stated as such: the passage from the
$\Lambda$-weighted count to unweighted twin counts is a
partial-summation argument whose error terms need writing out (higher
prime powers are disposed of by Lemma~\ref{l:c21}); and the
\emph{zero-oscillation caveat}---the drift law treats the oscillatory
explicit-formula contributions to class differences as having
vanishing logarithmic mean, the exact analogue of the
Rubinstein--Sarnak GRH-plus-linear-independence framework
\cite{RubinsteinSarnak}, and a conspiracy among hypothetical zeros
could in principle mimic or cancel part of the drift; clause (i) is
conjectured on the standard side of that hypothesis.

The mechanism is the twin analogue of the classical explanation of
Chebyshev's bias: prime-square terms contaminate the classes that can
contain them.  (Biases of twin primes in residue classes have been
studied before---Sahoo \cite{Sahoo} for the mod-$4$ race, and
\cite{PairBias} for prime pairs versus isolated primes---but we could
find no prior statement of the mod-$5$ race, the $(q^2-2,q^2)$
mechanism, or the quantified surplus below.)  Mod $3$ kills the pattern
$(q^2,q^2+2)$ outright, an amusing local accident that halves the
mechanism and (as a mod-$3$ computation showed, overturning a first guess of a
quadratic-residue bias against class $4$)
concentrates the surplus entirely on class $1$.  Data to $4\cdot10^9$:
class counts $(3980017,\ 3982505,\ 3981922)$, $D_1$ within one noise
unit of $T$ throughout, no persistent leader, and the
control race $\pi_{t,2}-\pi_{t,4}$ wandering like a fair coin.
Mod-8 data to $10^9$: class counts
$(856684,\ 855807,\ 856046,\ 855967)$ for $a=1,3,5,7$;
$D_7=+212$ at $10^9$ against predicted $T=+254$ with noise ${\sim}1068$,
positive through most of the range (log-density $0.775$), and the three
pairwise $\{1,3,5\}$ controls all within one noise unit of zero---the
two projections of the mechanism behave coherently, none of it provable
at present.

\begin{conjecture}[Least primes in progressions; (a) for part (i)
\cite{HLS,Wagstaff,Fiori,Leung}, (b) for part (ii)]\label{c22}
Let $p(a,q)$ be the least prime $\equiv a\pmod q$ and
$U(a,q)=\Li(p(a,q))/\varphi(q)$.  Then:
\emph{(i)} $U\Rightarrow\mathrm{Exp}(1)$ marginally as $q\to\infty$;
and, under the additional hypothesis that the class processes are
asymptotically jointly independent in the extreme-value sense (which
marginal convergence alone does not supply),
$\max_a U-H_{\varphi(q)}$ converges to a Gumbel law
($H_n=\sum_{k\le n}1/k$), the order-statistics form of Wagstaff's
$\max_a p(a,q)\sim\varphi(q)\log^2q$;
\emph{(ii)} the finite-$q$ deficit, averaged over all reduced classes
$a$ at fixed $q$, decomposes at first order $1/\log q$ as
\[
1-\mathbb E_a\,[U]
\;=\;\frac{\theta_{\mathrm{disc}}(q)+\theta_{\mathrm{corr}}(q)}{\log q}
+o\!\Bigl(\frac1{\log q}\Bigr),
\]
where $\theta_{\mathrm{disc}}(q)$ is now \emph{defined by formula}
from a named control (the parity-aware Cram\'er--Bernoulli model:
candidates $a,\,a+q,\,a+2q,\dots$ are independently ``prime'' with
hazard $h_i=(q/\varphi(q))\cdot s_i/\log(a+iq)$, $s_i=2$ for odd
candidates and $0$ for even, matching density and parity;
$\theta_{\mathrm{disc}}(q)=(1-\mathbb E_a[U^{\mathrm{model}}])\log q$
is computable exactly for each $q$), and
$\theta_{\mathrm{corr}}(q)=(1-\mathbb E_a[U])\log q
-\theta_{\mathrm{disc}}(q)$ is the ordering term---the part not
explained by discreteness.  The conjecture is that
$\theta_{\mathrm{corr}}(q)>0$ persists as $q\to\infty$ \emph{along
prime moduli}, where smooth-modulus singular-series structure is
absent.
\end{conjecture}

The stratified experiment that clause (ii) now rests on---demanded in
substance by the fourth external review---separates moduli by
factorization type in the matched range $q\in[1500,6000]$: over $40$
\emph{prime} moduli, $\theta_{\mathrm{disc}}=0.847$ (essentially
constant across $q$) and $\theta_{\mathrm{corr}}=0.824\pm0.009$;
over $40$ \emph{$7$-smooth} moduli in the same range,
$\theta_{\mathrm{disc}}=4.32\pm0.22$ (the large $q/\varphi(q)$
hazards make discreteness dominant) while
$\theta_{\mathrm{corr}}=-2.32\pm0.22$ reverses sign.  The comparison
with Leung \cite{Leung}, raised by the third review, is thereby given
its answer-shaped experiment: Leung derives the exponential law under
a uniform Hardy--Littlewood hypothesis and reports discrepancies for
\emph{smooth} moduli, and our smooth stratum indeed behaves
anomalously---but the positive ordering term survives, with small
error bars, exactly where smooth-modulus effects cannot reach.
Whether $\theta_{\mathrm{corr}}$ tends to a constant, and the exact
functional form in $1/\log q$, remain open: data at $q\le6000$ cannot
distinguish $1/\log q$ from $(\log\log q)/\log q$-type
refinements---the aggregate $\theta\approx1.5$--$1.9$ over
$q\in[10^2,6\times10^3]$, with $\mathbb E_a[U]$ bottoming near
$0.737$ around $q\approx200$ and recovering through $0.776$ by
$q\approx6000$, is calibration, not a determination of the form.  The
averaging family is now fixed as: all moduli of the stated type in
dyadic ranges of $q$, per stratum.

\begin{remark}\label{r:c22}
The deficit is reproducible and is \emph{not} an artifact of the
time-change: a Cram\'er pseudo-prime control (independent Bernoulli
``primes'' with the correct densities and parity) shows roughly half
of it---the discreteness of large per-candidate hazards---while the
other half vanishes under any reshuffling of the actual prime residues
(iid labels, or a random permutation of the true residue sequence,
both restore $\mathbb E[U]=1.00$).  The ordered sequence of primes
fills residue classes measurably faster than any exchangeable model,
and the stratified experiment above shows the effect is not a
smooth-modulus artifact: it is largest and cleanest on prime moduli.
We measure $\theta_{\mathrm{corr}}$; we cannot yet derive it.  It is
the deepest unexplained number produced by this exercise.
\end{remark}

\begin{conjecture}[Fermat quotients and the Wieferich ledger; (a), cf.\
\cite{CDP}]\label{c23}
With $q_p=(2^{p-1}-1)/p\bmod p$, three logically independent clauses:
\emph{(i)} [global equidistribution] $q_p/p$ equidistributes on
$[0,1)$ with discrepancy at the calibrated random-model scale:
$D_{\mathrm{KS}}(x)=O\bigl(\sqrt{\log\log\pi(x)/\pi(x)}\bigr)$
(the Chung--Smirnov law-of-the-iterated-logarithm scale; an earlier
draft demanded $O(\pi(x)^{-1/2})$, which is \emph{stronger than the
random model supports}---the same over-demand that killed the Mertens
conjecture, caught here by external review), with no drift;
\emph{(ii)} [shrinking targets] for every fixed $K\ge1$, and uniformly
for $K=K(x)\le\log x$,
$\#\{p\le x:\ q_p<K\}\sim\sum_{p\le x}\min(1,K/p)$---the fixed-$K$
case down to $K=1$ is what governs the zero event and does NOT follow
from (i), since $\{q_p=0\}$ is a target of shrinking measure $1/p$;
\emph{(iii)} [Wieferich count, fluctuation-calibrated]
$W(x)=\#\{\text{Wieferich } p\le x\}
=\log\log x+O_\varepsilon\bigl((\log\log x)^{1/2+\varepsilon}\bigr)$,
the fluctuation scale being that of the independent-Bernoulli model
(variance $\sim\log\log x$).
\end{conjecture}

Two corrections from external review are absorbed here.  First, an
earlier draft derived (iii) from (i) with ``hence''; that inference is
invalid---a square-root-size error in the aggregate distribution is
astronomically larger than $\log\log x$, and only the shrinking-target
clause (ii) at bounded $K$ speaks to the zero event.  Second, the
earlier fluctuation claim $O(1)$ violated this paper's own calibration
principle (criterion of demanding exactly the fluctuations the model
earns): with event probabilities $1/p$ the count has standard deviation
$\sim\sqrt{\log\log x}$, and the statement now says so.  We also
retract any location prediction for a third Wieferich prime: the
expected count reaching $3$ ``at doubly exponential heights'' is an
expected-count statement whose additive constant is uncalibrated, not a
forecast.  Consistency data: $\{1093,3511\}$ below $10^8$ against
expected $\approx2.7$; $\sqrt n\,\mathrm{KS}=0.82,\ 0.82,\ 0.52,\
0.63$ at $x=10^5,\dots,10^8$ (bounded, no drift); census at $K=100$:
observed $169$ against model $161.2$.

\begin{conjecture}[Conjecture F as a family; (a), \cite{HL1923},
cf.~\cite{JW}]\label{c24}
For odd $A$ let $Q_A(N)=\#\{n\le N:\ n^2+n+A \text{ prime}\}$ and
$C(A)=C(x^2+x+A)$.  Then, for every fixed $B>0$,
$Q_A(N)=C(A)\,I_A(N)\,(1+o(1))$ uniformly over odd
$1\le A\le(\log N)^B$, with residuals
$z_A(N)=(Q_A-C(A)I_A)/\sqrt{C(A)I_A}$ obeying the analogue of
Conjecture~\ref{c16}(ii) with the kernel now derived: two members
share the index set $n\le N$, so
\[
\operatorname{Cov}(Q_A,Q_{A'})\;\sim\;
\bigl[C(A,A')-C(A)\,C(A')\bigr]\,I_{A,A'}(N),
\qquad
C(A,A')=C(x^2+x+A,\ x^2+x+A'),
\]
with $C(A,A')=0$ (hence \emph{negative} correlation) for the locally
exclusive pairs, and the residual field converges to a centered
Gaussian field with the corresponding correlation kernel.
(Uniformity over a range growing with $N$ is the substantive claim; as
an external review noted, uniformity over any \emph{fixed} finite set
of $A$ is logically automatic from the individual asymptotics.)  Among
odd $A\le199$, Euler's $A=41$ has the largest constant,
$C(41)=6.6395463\ldots$ (Cohen's high-precision value \cite{Cohen}).
\end{conjecture}

Numerical evaluation of the kernel (all odd $A,A'\le99$ at
$N=10^6$: $1{,}225$ pairs, of which $289$ locally exclusive) returned
a structurally clean answer: the positive correlations of admissible
pairs and the negative correlations of exclusive pairs \emph{cancel},
mean off-diagonal $\rho=-0.003$, so the derived cross-member kernel is
null and the members are asymptotically independent after centering.
The observed profile variance ${\approx}0.37$ over the same range is
therefore \emph{not} a common-mode effect; it must be diagonal---each
$\operatorname{Var}(Q_A)$ sub-Poisson through within-sequence pair
correlations, the quadratic-family analogue of the
Montgomery--Soundararajan deficit (Conjecture~\ref{c20})---and part of
the spread is truncation noise in the constants themselves (the
$\pm10^{-3}$ wobble moves each $z_A$ by ${\sim}0.3$).  Deriving the
diagonal law is the family's remaining open component; the
cross-member question the review posed is now closed by computation.
At $N=10^6$: mean ratio $1.0006$, standard deviation $0.0027$,
$\max|z|=1.92$, rank correlation $0.99988$; for $A=41$, observed
$261{,}080$ against predicted $261{,}017.6$ ($z=+0.12$, using Cohen's
constant).  A cautionary note from the referee audit belongs here: our
own truncated Euler product for $C(41)$ at cutoff $2\cdot10^5$ reads
$6.64092$---wrong in the fourth significant decimal---because these
conditionally convergent products drift at the $10^{-3}$ scale over
practical cutoffs.  Quoted digits for all quadratic-family constants in
this paper are limited by the stated truncation wobble, not by the last
printed digit.  The entire
$10^8$-bit primality dataset compresses to one product formula per
$A$---the description-length virtue in its purest form.

\begin{conjecture}[The Goldbach lane race; apparently new]\label{c25}
For $n\equiv2\pmod4$ let $R_1(n)$, $R_3(n)$ be the ordered Goldbach
representation counts in the lanes $p\equiv q\equiv1$ and
$p\equiv q\equiv3\pmod4$ (Conjecture~\ref{c15}), and
$D(n)=R_3(n)-R_1(n)$.  Prime-square terms $n=q^2+p$ ($q$ odd prime)
land \emph{exclusively} in the $(1,1)$ lane, since
$q^2\equiv1\pmod4$ and then $n-q^2\equiv1\pmod4$ automatically.
Consequently the $(3,3)$ lane leads on average: with
\[
D_{\mathrm{sys}}(n)\;=\;\frac{2}{\bar\ell(n)}
\sum_{\substack{q \text{ odd prime},\ q\le\sqrt n\\ n-q^2 \text{ prime}}}
\log q\,\log(n-q^2),
\qquad
\bar\ell(n)\;=\;\frac{n-6}{\displaystyle\int_3^{n-3}
\frac{dt}{\log t\,\log(n-t)}},
\]
(the analytic mean of $\log p\log(n-p)$ under the lane profile; the
empirical lane mean is its estimator).  Write $\mathbb E_x$ for the
average over $n$ drawn from $\{n\equiv2\ (4),\ n\le x\}$ with weight
proportional to $1/n$ (the logarithmic sampling measure, now
explicit).  Then:
\emph{(i)} [drift law]
$\mathbb E_x[D-D_{\mathrm{sys}}]=o(\mathbb E_x[D_{\mathrm{sys}}])$ as
$x\to\infty$: the $(3,3)$ lane leads on average by exactly the square
contamination;
\emph{(ii)} [weighted mean and internal null lane]
$\mathbb E_x[D_{\mathrm{sys}}]$ is given by the Hardy--Littlewood
(Conjecture H) local model for $n-q^2$, whose local factors vary with
$n$ as the fourth external review stressed---most drastically at
$p=3$: for $n\equiv1\pmod3$, $3\mid n-q^2$ for \emph{every} prime
$q>3$, so the contamination collapses to the single $q=3$ term and
the race carries essentially no drift on the subprogression
$n\equiv10\pmod{12}$, while the drift concentrates on
$n\not\equiv1\pmod3$---an internal null lane, predicted by the same
mechanism that predicts the lead;
\emph{(iii)} [sign density] the per-$n$ drift-to-noise ratio
$\kappa(n)=D_{\mathrm{sys}}(n)\big/\sqrt{\SG(n)J(n)}$
(where $J(n)=\int_3^{n-3}dt/(\log t\log(n-t))$, so that
$\SG(n)J(n)$ is the total ordered representation count) satisfies
$\kappa(n)\asymp c(n)/\log n\to0$; under the Gaussian fluctuation
model the logarithmic density of $\{D>0\}$ is therefore $\tfrac12$ in
the limit, with the finite-$x$ sign fraction predicted to be
$\mathbb E_x[\Phi(\kappa)]$ ($\Phi$ the standard normal distribution
function)---decaying, but far from $\tfrac12$ at any accessible
height.
\end{conjecture}

Derivation status: the identity behind $D_{\mathrm{sys}}$ is the
$\Lambda$-weighted lane symmetry minus its prime-power part; the
transfer to unweighted counts is a partial-summation argument whose
error terms (higher prime powers, which enter only at the $n^{1/3}$
scale, and the variation of the weight across the lane) remain to be
written out---the same open transfer flagged at
Conjecture~\ref{c21}---and clause (iii) inherits the zero-oscillation
caveat stated there.

This is the Goldbach-partition analogue of Chebyshev's bias, produced
by the same explicit-formula mechanism as Conjecture~\ref{c21} but
through $q^2+p$ patterns rather than twin patterns; the single-sign
prediction (every square contamination falls into one lane) makes it
cleaner than the classical race.  We could find no prior statement.
It replaces, at the suggestion of an external review, a restatement of
the Lemke Oliver--Soundararajan mod-3 law that formerly occupied this
slot and carried no new content (their conjecture remains cited at
Conjecture~\ref{c16}'s family level and its verification data---%
$\hat c$ drifting $0.400\to0.370$ over $10^7\!\to\!4\cdot10^9$ with
$(1,1)/(2,2)$ symmetry $0.99989$---remain in the repository).

Verification.  Over $500$ log-spaced samples $n\le10^8$: mean
$D=76.5$ against mean empirical $D_{\mathrm{sys}}=64.5$, the
$(3,3)$-lane lead significant at $5.2$ standard errors of the mean.
Clause (ii)'s weighted model, computed in presieve-exact form
(trial division of each $n-q^2$ to $10^3$ with the exact Mertens
normalization), reproduces the empirical contamination with no fitted
parameter: model mean $66.3$ against empirical $67.8$ (ratio
$0.98$) on a $400$-sample run.  The internal null lane behaves as
predicted: on $n\equiv1\ (3)$ ($123$ samples) the model and empirical
$D_{\mathrm{sys}}$ both vanish and the measured lead drops to
$31\pm23$ (consistent with zero), while on $n\not\equiv1\ (3)$
($277$ samples) the lead is $88\pm21$ against predicted
contamination $96$---the drift lives exactly where the mechanism puts
it.  Clause (iii): computed $\mathbb E_x[\Phi(\kappa)]=0.572$
($\kappa$ ranging $0$--$0.64$ with mean $0.18$) against observed sign
fraction $0.588$; weak, and quantitatively the weakness the model
requires.

\section{The adversarial audit}\label{sec:audit}

Each statement above was attacked three ways.

\emph{Layer 1: literature.}  Five independent literature searches
combed OEIS, arXiv, and the standard references for prior art, producing the
novelty labels and the attribution section below.  This layer also
falsified one historical claim in our working notes (the
$k\ge1$ Stern list of Conjecture~\ref{c14} was described as partly new;
it is OEIS A060003 verbatim, and $1493$ is the largest known Stern
prime); the notes were corrected on the record.  The layer also
produced the audit programme's own worst failure: a snippet-depth
search cleared the primorial-twin statement that Lillie \cite{Lillie}
had in fact already made (Finding~\ref{f:c10}).  All novelty searches
for the statements admitted in the present revision (the
Conjecture~\ref{c1} and~\ref{c5} family lifts, the chain of
Conjecture~\ref{c6}, the null race of Conjecture~\ref{c8}, the
factorial twins of Conjecture~\ref{c10}) were therefore run at
abstract depth: the nearest-neighbour papers were read, not skimmed
from search results.

\emph{Layer 2: scripted refutation.}  Every
falsifiable-by-instance statement was pushed at least a decade past its
original bound: exception hunts across $(10^8,10^9]$ for
Conjectures~\ref{c13}, \ref{c14}, \ref{c15}, \ref{c17}; the uniformity
of Conjecture~\ref{c16} stressed on $d\le6000$; the statistical laws
re-tested at $4\cdot10^9$ (Conjectures~\ref{c18}--\ref{c21},
\ref{c25}, and the quintuplet count of Conjecture~\ref{c4}); the
recovery clause of Conjecture~\ref{c22} tested on fresh moduli
$q\in(3000,6000]$.  Outcome: one kill (Finding~\ref{f:c13}); no other
statement moved.  The statements admitted or upgraded in the present
revision were verified at production scale before admission: the
Conjecture~\ref{c1} profile across $150$ shifts; the chain of
Conjecture~\ref{c6} over five decades to $10^7$; the null race of
Conjecture~\ref{c8} at $10^7$; factorial twins to $n=700$; the
CRT-exact factorization of Conjecture~\ref{c11} through a joint
period of $1.2\times10^8$ with the scan extended to $n=6000$; the
covariance kernels of Conjectures~\ref{c16}(ii) and \ref{c24}(ii)
evaluated over $870$ and $1{,}225$ pairs; the stratified experiment
of Conjecture~\ref{c22} on $80$ moduli; and the weighted drift and
internal null lane of Conjecture~\ref{c25} on $400$ fresh samples.

\emph{Layer 3: clean-context replication.}  Independent replications,
working from the bare statements alone without access to our code or
notes, recomputed constants from definitions and recounted sequences
with their own implementations.  All recomputed constants landed inside
the stated truncation error; all recounts agreed (one replication
initially disagreed at $z\approx-22$, traced by the replicator
itself to a bug in \emph{its} sieve---every prime $n\le Q$ had been
struck out by its own modulus---after which it agreed at $z<1.6$).  The
same layer reproduced the Fibonacci deficit of Remark after
Conjecture~\ref{c9} and the ordering anomaly of Remark~\ref{r:c22},
confirming both are properties of the primes rather than of our code.

We draw one methodological conclusion.  Every fatal mathematical
error the audit programme has caught in its own output was not
statistical but \emph{algebraic}: factorizations or congruence
collapses living on density-zero or positive-density-but-structured
subsequences (Findings~\ref{f:c13} and~\ref{f:c4}, the $k=4$ parity
branch of Conjecture~\ref{c5}, the inadmissible naive chain of
Conjecture~\ref{c6}); and the one fatal \emph{bibliographic} error
was a search read at snippet depth (Finding~\ref{f:c10}).
Probabilistic sanity checks, however extensive, integrate over
density and cannot see such families; only structural attack surfaces
(here: ``test the claimed exception census on special subsequences'')
find them.  An audit that only re-runs the original verification at a
larger bound would have declared the false version of
Conjecture~\ref{c13} sound.

\section{Attribution and novelty}\label{sec:attrib}

Previously stated in essence, with our contribution being the uniform
quantification, computed constants, and re-verification:
C\ref{c1}$_{d=2}$ (OEIS A080149; \cite{Carella});
C\ref{c3} \cite{Dubner2002};
C\ref{c6}'s length-2 sub-chain (constant catalogued as OEIS A188596);
C\ref{c9} \cite{GranthamGranville};
C\ref{c12} \cite{CaldwellGallot};
C\ref{c13}-core \cite{HL1923};
C\ref{c14} (OEIS A060003);
C\ref{c15} \cite{Martin};
C\ref{c16}-core (uniform Hardy--Littlewood B, cf.\ \cite{GL});
C\ref{c17} \cite{Dubner2000};
C\ref{c18} \cite{Kourbatov};
C\ref{c23} \cite{CDP};
C\ref{c24}-core \cite{HL1923,JW}.
Retired former occupants, previously stated: the primorial twins of
slot \ref{c10}, Lillie \cite{Lillie} (see Finding~\ref{f:c10}); the
primorial benchmark before that and the Lemke Oliver--Soundararajan
restatement of slot \ref{c25}, \cite{CaldwellGallot} and \cite{LOS}
(the third external review caught a stale attribution table for
these, corrected on the record); the Sophie Germain statement of slot
\ref{c6} and the quadratic triple of slot \ref{c8}, retained inside
the new statements as attributed calibration.

For the remainder we distinguish, following the external review, two
grades of novelty.  \emph{Bibliographic novelty only}---unrecorded
specializations of Bateman--Horn, analogous to evaluating a classical
special function at a new argument, offered as benchmark data: the
constants of C\ref{c2}, C\ref{c5}(iii) and C\ref{c7}.
\emph{Conceptual novelty}---a mechanism, law, family uniformity, or
dichotomy we could not find in the literature:
C\ref{c1}'s uniform quadratic de Polignac family (the profile law,
not the $d=2$ instance);
C\ref{c4}'s power-obstruction ladder;
C\ref{c5}'s complete twin-base family analysis;
C\ref{c6}'s alternating cyclotomic chain;
C\ref{c8}'s null-mechanism contrast prediction;
C\ref{c10}'s factorial-twin finiteness (the convergent
Borel--Cantelli exemplar, admitted after the deeper search of
Finding~\ref{f:c10});
C\ref{c11}'s entanglement-aware local product for a mixed
polynomial--exponential sequence;
C\ref{c16}(ii)'s derived covariance kernel (and its C\ref{c24}(ii)
family analogue with the computed null cross-kernel);
C\ref{c19}'s liminf constant
(against the slope-1 laws of Shanks, Wolf, and Kourbatov--Wolf
\cite{KW2020}); C\ref{c20}'s microscopic variance law (a sharpening
of \cite{MS}, with the mesoscopic numerics of \cite{SL} as
precedent); C\ref{c21}'s race mechanism (against the twin-bias
literature \cite{Sahoo,PairBias}); C\ref{c22}(ii)'s deficit
decomposition with its stratified differentiating experiment; and
C\ref{c25}'s Goldbach lane race.  Among these we
consider C\ref{c21}, C\ref{c25} and C\ref{c8} (the
square-contamination mechanism family with its positive and negative
controls), C\ref{c22}(ii) (the least-prime deficit law with its
ordering anomaly), and the derived kernels of C\ref{c16}(ii) and
C\ref{c24}(ii) the paper's most genuinely new content, and
Question~\ref{q:theta} its best export.

\begin{question}\label{q:theta}
Determine $\theta(q)$ of Conjecture~\ref{c22}: derive, from the
Hardy--Littlewood correlations or otherwise, the first-order deficit of
$\mathbb E_a\,\Li(p(a,q))/\varphi(q)$ below its random-model value
$1$---equivalently, explain quantitatively why the primes in their
natural order occupy residue classes faster than any exchangeable
resampling of themselves.
\end{question}

\begin{thebibliography}{99}\small

\bibitem{HL1923} G.~H. Hardy and J.~E. Littlewood,
\emph{Some problems of `Partitio numerorum'; III: On the expression of
a number as a sum of primes}, Acta Math.\ 44 (1923), 1--70.
DOI: 10.1007/BF02403921.

\bibitem{BH} P.~T. Bateman and R.~A. Horn,
\emph{A heuristic asymptotic formula concerning the distribution of
prime numbers}, Math.\ Comp.\ 16 (1962), 363--367.
DOI: 10.1090/S0025-5718-1962-0148632-7.

\bibitem{HB} D.~R. Heath-Brown,
\emph{Primes represented by $x^3+2y^3$},
Acta Math.\ 186 (2001), 1--84.  DOI: 10.1007/BF02392715.

\bibitem{CaldwellGallot} C.~K. Caldwell and Y.~Gallot,
\emph{On the primality of $n!\pm1$ and
$2\times3\times5\times\cdots\times p\pm1$},
Math.\ Comp.\ 71 (2002), 441--448.
DOI: 10.1090/S0025-5718-01-01315-1.

\bibitem{Lillie} G.~Lillie,
\emph{About the primality of primorials},
arXiv:2110.04302 [math.NT] (2021).

\bibitem{RubinsteinSarnak} M.~Rubinstein and P.~Sarnak,
\emph{Chebyshev's bias},
Experiment.\ Math.\ 3 (1994), no.~3, 173--197.
DOI: 10.1080/10586458.1994.10504289.

\bibitem{GranthamGranville} J.~Grantham and A.~Granville,
\emph{Fibonacci primes, primes of the form $2^n-k$ and beyond},
J.~Number Theory 261 (2024), 190--219; arXiv:2307.07894.
DOI: 10.1016/j.jnt.2024.02.002.

\bibitem{Wagstaffheu} S.~S. Wagstaff, Jr.,
\emph{Divisors of Mersenne numbers},
Math.\ Comp.\ 40 (1983), 385--397.
DOI: 10.1090/S0025-5718-1983-0679454-X.

\bibitem{Dubner2000} H.~Dubner,
\emph{Twin prime conjectures},
J.~Recreational Math.\ 30, no.~3 (1999--2000) (article reproduced at
OEIS A007534).

\bibitem{Dubner2002} H.~Dubner,
\emph{Carmichael numbers of the form $(6m+1)(12m+1)(18m+1)$},
J.~Integer Seq.\ 5 (2002), Article 02.2.1.

\bibitem{Martin} K.~Martin,
\emph{Refined Goldbach conjectures with primes in progressions},
Exp.\ Math.\ 31 (2022), 226--232; arXiv:1806.00946.
DOI: 10.1080/10586458.2019.1596849.

\bibitem{GL} D.~A. Goldston and A.~H. Ledoan,
\emph{Jumping champions and gaps between consecutive primes},
Int.\ J. Number Theory 7 (2011), 1413--1421; arXiv:0910.2960.
DOI: 10.1142/S179304211100471X.

\bibitem{Kourbatov} A.~Kourbatov,
\emph{Maximal gaps between prime $k$-tuples: a statistical approach},
J.~Integer Seq.\ 16 (2013), Article 13.5.2; arXiv:1301.2242.  See also
arXiv:1309.4053.

\bibitem{Shanks} D.~Shanks,
\emph{On maximal gaps between successive primes},
Math.\ Comp.\ 18 (1964), 646--651.
DOI: 10.1090/S0025-5718-1964-0167472-8.

\bibitem{Granville} A.~Granville,
\emph{Harald Cram\'er and the distribution of prime numbers},
Scand.\ Actuar.\ J.\ 1995:1, 12--28.
DOI: 10.1080/03461238.1995.10413946.

\bibitem{Gallagher} P.~X. Gallagher,
\emph{On the distribution of primes in short intervals},
Mathematika 23 (1976), 4--9.  DOI: 10.1112/S0025579300016442.

\bibitem{MS} H.~L. Montgomery and K.~Soundararajan,
\emph{Primes in short intervals},
Comm.\ Math.\ Phys.\ 252 (2004), 589--617.
DOI: 10.1007/s00220-004-1222-4.

\bibitem{HLS} T.~Haddad, S.-K. Leung, and C.~Sabuncu,
\emph{Visiting early at prime times}, arXiv:2408.11781 (2024).
DOI: 10.48550/arXiv.2408.11781.

\bibitem{Wagstaff} S.~S. Wagstaff, Jr.,
\emph{Greatest of the least primes in arithmetic progressions having a
given modulus}, Math.\ Comp.\ 33 (1979), 1073--1080.
DOI: 10.1090/S0025-5718-1979-0528061-7.

\bibitem{CDP} R.~Crandall, K.~Dilcher, and C.~Pomerance,
\emph{A search for Wieferich and Wilson primes},
Math.\ Comp.\ 66 (1997), 433--449.
DOI: 10.1090/S0025-5718-97-00791-6.

\bibitem{JW} M.~J. Jacobson, Jr.\ and H.~C. Williams,
\emph{New quadratic polynomials with high densities of prime values},
Math.\ Comp.\ 72 (2003), 499--519.
DOI: 10.1090/S0025-5718-02-01418-7.

\bibitem{LOS} R.~J. Lemke Oliver and K.~Soundararajan,
\emph{Unexpected biases in the distribution of consecutive primes},
Proc.\ Natl.\ Acad.\ Sci.\ USA 113 (2016), E4446--E4454;
arXiv:1603.03720.  DOI: 10.1073/pnas.1605366113.

\bibitem{KW2020} A.~Kourbatov and M.~Wolf,
\emph{On the first occurrences of gaps between primes in a residue
class}, J.~Integer Seq.\ 23 (2020), Article 20.9.3; arXiv:2002.02115.

\bibitem{SL} M.-A. Sanchis-Lozano,
\emph{A heuristic study of the distribution of primes in short and
not-so-short intervals}, arXiv:1804.07659 (2018).
DOI: 10.48550/arXiv.1804.07659.

\bibitem{Sahoo} S.~Sahoo,
\emph{On twin prime distribution and associated biases},
arXiv:2111.09053 (2021).  DOI: 10.48550/arXiv.2111.09053.

\bibitem{PairBias} W.~Puszkarz,
\emph{Statistical bias in the distribution of prime pairs and isolated
primes}, arXiv:1807.00406 (2018).  DOI: 10.48550/arXiv.1807.00406.

\bibitem{Leung} S.-K. Leung,
\emph{Moments of primes in progressions to a large modulus},
arXiv:2402.07941.  DOI: 10.48550/arXiv.2402.07941.

\bibitem{Fiori} A.~Fiori,
\emph{The least prime in arithmetic an progression} [sic],
arXiv:2404.02329 (2024).  DOI: 10.48550/arXiv.2404.02329.

\bibitem{Cohen} H.~Cohen,
\emph{High precision computation of Hardy--Littlewood constants},
unpublished notes (PDF hosted at OEIS A221712).

\bibitem{Carella} N.~A. Carella,
\emph{Twin primes in quadratic arithmetic progressions},
arXiv:1710.07827 (unrefereed preprint).
DOI: 10.48550/arXiv.1710.07827.

\end{thebibliography}

\end{document}
